\documentclass[11pt]{article}
\usepackage[utf8]{inputenc}
\usepackage[margin=1in]{geometry}
\usepackage{amsmath}
\usepackage{graphicx}
\usepackage{booktabs}
\usepackage{hyperref}
\usepackage[round]{natbib}
\usepackage{lineno}
\usepackage{xcolor}
\usepackage{multirow}
\usepackage{float}

\linenumbers

\title{Glutamate Neurotransmission from Leptin Receptor Cells Is Required for Typical Puberty and Reproductive Function in Female Mice}

\author{
Author A$^{1,*}$, Author B$^{1}$, Author C$^{1}$, Author D$^{2}$, Author E$^{1}$ \\
\\
$^{1}$Department of Molecular and Integrative Physiology, University of Michigan \\
$^{2}$Department of Internal Medicine, University of Michigan \\
$^{*}$Corresponding author: authora@umich.edu
}

\date{}

\begin{document}
\maketitle

\begin{abstract}
Leptin is an essential metabolic signal for the reproductive axis, yet the neural circuit mechanisms linking leptin to gonadotropin-releasing hormone (GnRH) secretion remain incompletely understood. The ventral premammillary nucleus (PMv) is a glutamatergic hypothalamic nucleus densely expressing leptin receptor (LepRb) that projects directly to GnRH neurons and kisspeptin populations. Previous work concluded that GABAergic, not glutamatergic, leptin signaling mediates reproductive function, based on conditional Vglut2 deletion across all LepRb neurons. Here we demonstrate using three complementary approaches that glutamate neurotransmission from LepRb PMv neurons is required for leptin's control of puberty and fertility. First, chemogenetic (DREADD) activation of LepRb PMv neurons induced transient LH release in fed adult females (15 of 18 mice showed successful PMv targeting; peak LH = $2.14 \pm 0.42$~ng/ml vs. baseline $0.38 \pm 0.09$~ng/ml, $p < 0.001$, paired $t$-test), with LH amplitude correlated with cFos-immunoreactive PMv neuron number ($r = 0.72$, $p < 0.01$). Second, conditional Vglut2 deletion in LepRb neurons (LepRb-$\Delta$Vglut2) delayed first estrus ($t_{14} = 2.45$, $p = 0.028$), disrupted estrous cycles, and increased GnRH content in median eminence terminals ($t_5 = 8.01$, $p = 0.0005$), indicating impaired GnRH release. Third, in a PMv-specific rescue model, restoring LepRb expression in the PMv rescued vaginal opening, follicle maturation, and pregnancy, while simultaneous Vglut2 deletion abolished pubertal development entirely. These results establish that glutamatergic neurotransmission from PMv LepRb neurons is essential for leptin's reproductive actions.
\end{abstract}

\section{Introduction}

Reproduction is an energetically costly process that requires adequate metabolic reserves. Leptin, an adipocyte-derived hormone signaling energy sufficiency, is a critical permissive factor for the reproductive axis \citep{ahima1996leptin, barash1996leptin}. Leptin deficiency in both humans and mice results in hypogonadotropic hypogonadism and infertility, which are rescued by leptin administration. However, the neural circuits through which leptin communicates metabolic status to the GnRH pulse generator remain incompletely characterized \citep{herbison2016gnrh}.

The ventral premammillary nucleus (PMv) of the hypothalamus is a glutamatergic nucleus in which approximately 90\% of neurons express the leptin receptor (LepRb) \citep{scott2009leptin, donato2011leptin}. PMv neurons project directly to GnRH cell bodies and to kisspeptin-expressing neurons in both the arcuate nucleus and the anteroventral periventricular nucleus (AVPV) \citep{ross2018pmv}. Despite this strategic positioning, the PMv has been relatively overlooked in models of leptin-mediated reproductive control.

A key finding that shaped the field was that conditional deletion of LepRb from all Vglut2-expressing neurons (Vglut2-Cre $\times$ LepRb-floxed) caused virtually no reproductive deficits, while deletion from GABAergic neurons produced severe impairment \citep{vong2011leptin, zuure2013leptin}. This led to the prevailing view that GABAergic leptin signaling is the primary mediator of reproductive function. However, the global Vglut2-neuron deletion removes leptin signaling from diverse glutamatergic populations throughout the brain, potentially engaging compensatory mechanisms that mask the specific contribution of PMv neurons.

Here we directly test the role of PMv glutamate signaling using three approaches: chemogenetic activation, conditional Vglut2 deletion in LepRb neurons, and PMv-specific LepRb rescue with and without Vglut2.

\section{Results}

\subsection{Chemogenetic Activation of PMv LepRb Neurons Induces LH Release}

LepRb-Cre mice received unilateral stereotaxic injections of AAV-hM3Dq into the PMv. After recovery, serial blood sampling was performed via the Culex Automated Blood Sampling System before and after intravenous CNO (0.5~mg/kg) (Table~\ref{tab:dreadd}).

\begin{table}[H]
\centering
\caption{LH responses following chemogenetic activation of PMv LepRb neurons. PMv-hits classified by mCherry$^+$ cFos$^+$ neurons in PMv post-hoc. Baseline LH is mean of 3 pre-injection samples (10~min intervals). Peak LH is maximum post-CNO value within 60~min. Statistical analysis by paired $t$-test (baseline vs. peak) or Kruskal--Wallis test across groups.}
\label{tab:dreadd}
\begin{tabular}{lcccccc}
\toprule
\textbf{Group} & \textbf{$n$} & \textbf{Baseline LH (ng/ml)} & \textbf{Peak LH (ng/ml)} & \textbf{$\Delta$LH} & \textbf{$t$ or $H$} & \textbf{$p$-value} \\
\midrule
PMv-hits (CNO) & 15 & $0.38 \pm 0.09$ & $2.14 \pm 0.42$ & $+1.76$ & $t = 4.21$ & $< 0.001$ \\
PMv-misses (CNO) & 2 & $0.41 \pm 0.12$ & $0.48 \pm 0.14$ & $+0.07$ & --- & NS \\
mCherry ctrl (CNO) & 5 & $0.36 \pm 0.08$ & $0.39 \pm 0.10$ & $+0.03$ & --- & NS \\
PMv-hits (Cloz) & 4 & $0.42 \pm 0.11$ & $1.87 \pm 0.38$ & $+1.45$ & $t = 3.84$ & $< 0.05$ \\
Clozapine alone & 7 & $0.35 \pm 0.07$ & $0.38 \pm 0.09$ & $+0.03$ & --- & NS \\
\midrule
\multicolumn{5}{l}{LH vs. cFos$^+$ PMv neurons (PMv-hits):} & $r = 0.72$ & $p < 0.01$ \\
\bottomrule
\end{tabular}
\end{table}

Analysis of arcuate nucleus contamination revealed that some PMv-hit animals with LH increase showed AAV spread to the posterior Arc. cFos/mCherry colocalization in the posterior Arc differed significantly across groups (Kruskal--Wallis $H = 14.63$, $p < 0.0001$; Table~\ref{tab:arc}).

\begin{table}[H]
\centering
\caption{cFos/mCherry colocalization in the posterior arcuate nucleus across groups. Values are mean $\pm$ SEM per section. Kruskal--Wallis test with Dunn's post-hoc comparisons.}
\label{tab:arc}
\begin{tabular}{lccccc}
\toprule
\textbf{Group} & \textbf{$n$} & \textbf{cFos$^+$/mCherry$^+$ cells} & \textbf{vs. LH increase} & \textbf{vs. controls} & \textbf{$p$ (post-hoc)} \\
\midrule
PMv-hit, LH increase & 8 & $12.4 \pm 3.2$ & --- & $< 0.001$ & --- \\
PMv-hit, no LH increase & 7 & $3.1 \pm 1.4$ & $< 0.001$ & $< 0.05$ & --- \\
Negative controls & 7 & $0.8 \pm 0.4$ & $< 0.001$ & --- & --- \\
\bottomrule
\end{tabular}
\end{table}

\subsection{LepRb-$\Delta$Vglut2 Females Show Delayed Puberty and Disrupted Cyclicity}

Conditional Vglut2 deletion in LepRb neurons produced delayed first estrus and disrupted reproductive function selectively in females (Table~\ref{tab:puberty}).

\begin{table}[H]
\centering
\caption{Pubertal development in LepRb-$\Delta$Vglut2 mice. VO = vaginal opening; BPS = balanopreputial separation. Statistical analysis by Student's $t$-test. Values are mean $\pm$ SEM.}
\label{tab:puberty}
\begin{tabular}{llccccc}
\toprule
\textbf{Sex} & \textbf{Metric} & \textbf{Control} & \textbf{LepRb-$\Delta$Vglut2} & \textbf{$t$-stat} & \textbf{df} & \textbf{$p$-value} \\
\midrule
Female & Age of VO (days) & $28.4 \pm 0.8$ & $29.6 \pm 1.1$ & $0.84$ & 18 & $0.41$ \\
Female & Body mass at VO (g) & $14.2 \pm 0.6$ & $15.4 \pm 0.7$ & $1.47$ & 18 & $0.16$ \\
Female & Age first estrus (days) & $32.1 \pm 0.9$ & $36.8 \pm 1.6$ & $2.45$ & 14 & $\mathbf{0.028}$ \\
Female & Body mass at 1st estrus (g) & $16.8 \pm 0.7$ & $17.2 \pm 0.8$ & $0.59$ & 14 & $0.56$ \\
Male & Age of BPS (days) & $26.2 \pm 0.7$ & $28.4 \pm 1.0$ & $1.80$ & 11 & $0.10$ \\
Male & Body mass at BPS (g) & $18.1 \pm 0.6$ & $20.8 \pm 0.9$ & $2.43$ & 11 & $\mathbf{0.033}$ \\
\bottomrule
\end{tabular}
\end{table}

GnRH immunoreactivity in median eminence terminals was markedly increased in LepRb-$\Delta$Vglut2 females, indicating impaired GnRH secretion (Table~\ref{tab:gnrh}).

\begin{table}[H]
\centering
\caption{GnRH immunoreactivity (integrated density) in arcuate nucleus/median eminence. Values normalized to control mean. Statistical analysis by Student's $t$-test.}
\label{tab:gnrh}
\begin{tabular}{lccccc}
\toprule
\textbf{Region} & \textbf{Control ($n=4$)} & \textbf{LepRb-$\Delta$Vglut2 ($n=3$)} & \textbf{$t$-stat} & \textbf{df} & \textbf{$p$-value} \\
\midrule
Arc/ME GnRH-ir & $1.00 \pm 0.06$ & $2.84 \pm 0.18$ & $8.01$ & 5 & $\mathbf{0.0005}$ \\
\bottomrule
\end{tabular}
\end{table}

\subsection{PMv-Specific Rescue: Vglut2 Is Required for Pubertal Development}

In LepRloxTB mice (LepR-null, Cre-inducible), stereotaxic AAV-Cre injection into the PMv restored LepRb expression locally. When Vglut2 was simultaneously deleted (LepRloxTB;Vglut2$^{flox}$ + AAV-Cre), pubertal development was completely abolished (Table~\ref{tab:rescue}).

\begin{table}[H]
\centering
\caption{Reproductive outcomes in PMv-specific rescue experiments. pSTAT3-ir confirms LepRb restoration in PMv. Ovarian histology assessed at 12 weeks post-injection. Breeding test: 60-day cohabitation with proven males.}
\label{tab:rescue}
\begin{tabular}{lcccc}
\toprule
\textbf{Measure} & \textbf{WT} & \textbf{LepRloxTB + AAV-Cre} & \textbf{LepRloxTB;Vglut2$^{fl}$ + AAV-Cre} & \textbf{$p$-value} \\
\midrule
pSTAT3$^+$ in PMv & Yes & Restored & Restored & --- \\
Vglut2 mRNA in PMv & Normal & Normal & Decreased & $< 0.01$ \\
Vaginal opening & Yes & Yes & No & $< 0.001$ \\
Mature follicles & Present & Present & Absent & $< 0.001$ \\
Corpora lutea & Present & Present & Absent & $< 0.001$ \\
Estrous cycles & Normal & Restored & None & $< 0.001$ \\
Pregnancy (60-day test) & 4/4 & 3/4 & 0/4 & $< 0.01$ \\
\bottomrule
\end{tabular}
\end{table}

\subsection{Validation of Knockout Efficiency}

In situ hybridization confirmed elimination of Vglut2/Lepr colocalization in both VMHdm and PMv of LepRb-$\Delta$Vglut2 mice (Table~\ref{tab:ish}).

\begin{table}[H]
\centering
\caption{Vglut2 and Lepr mRNA colocalization by dual-label in situ hybridization. Percentage of Lepr$^+$ neurons also expressing Vglut2 in VMHdm and PMv. Values are mean $\pm$ SEM, $n = 4$ mice per genotype. Statistical analysis by Student's $t$-test.}
\label{tab:ish}
\begin{tabular}{lcccc}
\toprule
\textbf{Region} & \textbf{Control (\% coloc.)} & \textbf{LepRb-$\Delta$Vglut2 (\% coloc.)} & \textbf{$t$-stat} & \textbf{$p$-value} \\
\midrule
VMHdm & $78.4 \pm 4.2$ & $4.1 \pm 1.8$ & $16.2$ & $< 0.0001$ \\
PMv & $82.1 \pm 3.8$ & $3.6 \pm 1.5$ & $18.4$ & $< 0.0001$ \\
\bottomrule
\end{tabular}
\end{table}

\section{Discussion}

Our findings establish that glutamatergic neurotransmission from LepRb-expressing PMv neurons is essential for leptin's control of puberty and reproductive function in female mice. This conclusion is supported by three independent experimental approaches: chemogenetic activation, conditional Vglut2 deletion, and PMv-specific rescue with concomitant Vglut2 deletion.

The chemogenetic experiments demonstrate that acute, unilateral activation of LepRb PMv neurons is sufficient to induce transient LH release in fed adult females. The correlation between LH amplitude and cFos$^+$ PMv neuron number supports a dose-response relationship between PMv activity and GnRH/LH output \citep{donato2011leptin, ross2018pmv}. The confirmation that clozapine (a CNO metabolite) also induced LH release in hM3Dq-expressing animals validates the DREADD approach \citep{alexander2009remote}.

The LepRb-$\Delta$Vglut2 model reveals that glutamate release from LepRb neurons is required for normal pubertal timing and estrous cyclicity. The delayed first estrus and increased GnRH immunoreactivity in median eminence terminals---indicating accumulation due to impaired release---parallel the phenotype seen in leptin-deficient models \citep{barash1996leptin, herbison2016gnrh}. Notably, vaginal opening was not significantly delayed, suggesting that this earlier pubertal milestone is less dependent on glutamatergic LepRb signaling.

The most compelling evidence comes from the PMv-specific rescue experiment. Restoring LepRb expression solely in the PMv was sufficient to rescue vaginal opening, follicle maturation, estrous cycles, and fertility in otherwise LepR-null mice. When Vglut2 was simultaneously deleted in the PMv, all reproductive milestones were abolished. This demonstrates that glutamate release from PMv LepRb neurons is not merely modulatory but is required for the reproductive effects of leptin \citep{donato2011leptin, bellefontaine2014pmv}.

These results challenge the prevailing view that GABAergic leptin signaling alone mediates reproductive function \citep{vong2011leptin, zuure2013leptin}. The discrepancy with previous global Vglut2-LepRb deletion studies likely reflects compensatory mechanisms or the dilution effect of removing Vglut2 across all glutamatergic LepRb populations.

\section{Methods}

\subsection{Animals}
All procedures were approved by the University of Michigan IACUC (protocol PRO00010420). Mouse lines included LepRb-Cre (JAX 032457), LepRloxTB (JAX 018989), Vglut2-floxed (JAX 012898), and Kiss1-hrGFP (JAX 023425). Animals were maintained on low-phytoestrogen diets (Envigo 2016/2019) under 12:12 light cycle.

\subsection{Stereotaxic Surgery}
PMv coordinates: AP $-5.4$~mm, ML $\pm 0.5$~mm, DV $-5.4$~mm from dura. Approximately 50~nl of AAV was injected via glass pipette and picospritzer. Recovery period: 1 month.

\subsection{Serial Blood Sampling}
Mice were cannulated (aortic artery for sampling, jugular vein for injection) and connected to a Culex Automated Blood Sampling System. Baseline: 3 samples at 10~min intervals. Post-injection: every 10~min for 60~min. LH measured by ELISA \citep{steyn2013development}.

\subsection{Reproductive Phenotyping}
Vaginal opening and first estrus monitored daily from P21. Estrous cyclicity monitored by daily vaginal lavage for 21 days. Breeding tests: 60-day cohabitation with proven males. Ovarian histology: serial sectioning with H\&E staining.

\subsection{Immunohistochemistry and In Situ Hybridization}
GnRH immunoreactivity quantified by fluorescence microscopy with integrated density measurements (ImageJ). Dual-label in situ hybridization for Lepr and Vglut2 mRNA performed using RNAscope probes.

\subsection{Statistical Analysis}
Paired $t$-tests compared baseline vs. post-injection LH. Unpaired Student's $t$-tests compared control vs. knockout groups. Kruskal--Wallis test with Dunn's post-hoc used for multi-group non-parametric comparisons. All data presented as mean $\pm$ SEM; $\alpha = 0.05$.

\section{Conclusion}

We demonstrate that glutamatergic neurotransmission from LepRb-expressing PMv neurons is required for leptin's control of puberty onset and reproductive function. PMv-specific LepRb restoration rescues fertility, but only when glutamate signaling is intact. These findings revise the model of leptin-mediated reproductive control to include an essential glutamatergic component operating through the PMv.

\bibliographystyle{plainnat}
\bibliography{references}

\end{document}
